\documentclass[11pt,a4paper]{article}
\usepackage[margin=1in]{geometry}
\usepackage[T1]{fontenc}
\usepackage{lmodern}
\usepackage{booktabs}
\usepackage{array}
\usepackage{xcolor}
\usepackage{hyperref}
\usepackage{enumitem}
\usepackage{listings}

\hypersetup{colorlinks=true,linkcolor=blue,urlcolor=blue,citecolor=blue}
\setlist{nosep,leftmargin=*}
\lstset{basicstyle=\ttfamily\small,breaklines=true,frame=single}
\setlength{\emergencystretch}{3em}
\sloppy

\title{TLA+ Models of Parsl\\\large Project Report}
\author{TLA+ Parsl Audit Project}
\date{October 2026}

\begin{document}
\maketitle
\begin{abstract}
This project develops a small executable TLA+ abstraction of Parsl for finding
high-probability protocol errors. The primary focus is communication across process and thread
boundaries; asynchronous interleavings and terminal-state handling are secondary. The result is a
finite, repeatable audit connecting source transitions to TLC counterexamples, candidate safety
protocols, and runtime probes against the installed Parsl implementation.
\end{abstract}

\section{Research question}
Distributed workflow systems can fail even when the user function is correct. A task can leave one
queue and fail to arrive at the next component; a late retry result can overwrite a newer result;
or a callback can mutate a registry while failure fan-out is traversing it. The central question is:
\emph{when a logical task crosses asynchronous boundaries, who owns it, and what must remain true
when the next boundary fails?}

The model separates the following identities:
\begin{center}
\fbox{\parbox{0.82\textwidth}{\centering
logical task $\rightarrow$ physical attempt $\rightarrow$ transport message $\rightarrow$
Future and monitoring state}}
\end{center}

This separation lets the model classify an attempt-0 result as stale after attempt 1 has become
current, rather than allowing it to resolve the logical Future a second time.

\section{Model and evidence method}
TLA+ describes states and transitions rather than implementation syntax. Each concrete source-backed
finding follows a Current/Fixed discipline:

\begin{enumerate}
  \item Current preserves the source-like transition and should produce a meaningful counterexample.
  \item Fixed expresses a candidate safe protocol: retain ownership, isolate malformed input,
        check attempt generation, or snapshot a collection before callbacks.
  \item A focused Python probe exercises the installed Parsl path when the boundary is executable.
  \item The source location, trace, candidate fix, and limitations are recorded in the bug ledger.
\end{enumerate}

The resulting evidence chain is:
\begin{center}
\fbox{source transition $\rightarrow$ Current TLC trace $\rightarrow$ runtime probe $\rightarrow$
Fixed invariant}
\end{center}

A passing Fixed run is evidence for finite constants and transitions. It is not an unbounded proof
about Python, a real network, or every deployment.

\section{Coverage}
The model families cover:
\begin{itemize}
  \item DAG dependencies, Futures, retries, cancellation, timeouts, and \texttt{join\_app};
  \item logical tasks and physical attempts, serializer identity, multipart frames, routes,
        acknowledgements, duplicates, and stale results;
  \item HTEX, Thread, Flux, MPI, Work Queue, TaskVine, and Radical-Pilot lifecycle paths;
  \item provider provisioning, scheduler status, cancellation, scale-in/scale-out, and cleanup;
  \item DataFuture readiness, partial files, retries, atomic publication, and transfer cleanup;
  \item monitoring queues, database bookkeeping, shutdown, persistence, heartbeat, and timeout behavior.
\end{itemize}

\begin{table}[h]
\centering
\caption{Current project inventory}
\begin{tabular}{lr}
\toprule
Artifact & Count \\
\midrule
TLA+ modules & 676 \\
TLC configuration files & 1,478 \\
Unittest methods & 761 \\
Runtime probe files & 490 \\
Foundational runtime entries & 486 \\
Foundational TLC configurations & 754 \\
Bug ledger records & 318 \\
Unique numeric BUG identifiers (approximately) & 317 \\
\bottomrule
\end{tabular}
\end{table}

\section{Representative findings}
\subsection{HTEX task dispatch send failure}
In the HTEX interchange path, a pending task is removed from the queue before it is sent to a
manager over ZeroMQ. If the send raises, the task is no longer pending, assigned, or terminalized.
The Future can remain pending indefinitely. The Current model violates a no-task-loss invariant;
the Fixed model retains ownership until successful transfer or emits an explicit terminal transport
failure. This is recorded as BUG-335.

Recent communication refinements cover malformed command ingress (BUG-336), command-reply send
failure (BUG-337), and heartbeat ACK send failure (BUG-338). A BUG-004 refinement covers synthetic
manager-loss results that fail during transport. The shared safety principle is that a failed
transport operation must preserve ownership for retry or make failure explicit and terminal.

\subsection{Stale result after retry}
If attempt 0 times out and attempt 1 succeeds, a delayed attempt-0 result must not overwrite the
current Future. The model requires both logical task identity and the current physical-attempt
generation. The stale result is consumed or rejected without changing terminal state.

\subsection{Callback mutation during failure fan-out}
Several failure paths iterate a live task dictionary while synchronously invoking
\texttt{Future.set\_exception}. A callback can remove an entry during iteration, causing an
exception and leaving unrelated Futures pending. The candidate safe protocol snapshots entries
before invoking callbacks.

\section{Ledger distribution}
The canonical ledger contains 318 records, including a refinement that reuses BUG-259, and about
317 unique numeric identifiers. The distribution below describes record assignment, not defect
prevalence.

\begin{table}[h]
\centering
\caption{Ledger record distribution}
\begin{tabular}{lr}
\toprule
Category & Records \\
\midrule
Providers and scheduler adapters & 92 \\
Executors and worker lifecycle & 78 \\
Serialization and ZMQ transport & 53 \\
Monitoring and database & 32 \\
File staging and transfer & 29 \\
Clock, heartbeat, and timeout & 22 \\
\texttt{join\_app} and memoization & 9 \\
Core dataflow and Future lifecycle & 3 \\
\bottomrule
\end{tabular}
\end{table}

These records are source-backed protocol risks reproduced in the current implementation or in a
tightly scoped test double. They are not automatically confirmed upstream defects.

\section{Verification and reproduction}
The latest complete regression passed all 754 foundational TLC configurations and all 486 runtime
entries. The normal gate checks Fixed/safe configurations. A separate recent-model runner checks
that intentional Current configurations still produce their expected counterexamples.

\begin{lstlisting}[language=bash]
PYTHONPATH=/tmp/parsl-source \\
  PYTHON_BIN=/tmp/parsl-venv/bin/python \\
  PARSL_SOURCE=/tmp/parsl-source \\
  bash scripts/runtime_foundational_smoke.sh

JAVA_BIN=/path/to/java \\
  TLA_JAR=/path/to/tla2tools.jar \\
  TLC_SIMULATE=100 \\
  bash scripts/tlc_foundational_smoke.sh
\end{lstlisting}

\section{Interpretation and limits}
The abstraction omits unbounded Python heaps, arbitrary object graphs, operating-system scheduling,
real network timing, and many backend-specific implementation details. A Current counterexample
demonstrates a reachable source-like behavior under the model; it does not imply that every
production deployment fails. A runtime probe demonstrates behavior in a controlled setup; it is
not a production-scale experiment. The ledger is a research record, not an upstream issue tracker.

The project stops expanding when the selected communication paths and asynchronous races have
source-backed models, runtime evidence, documentation, and passing smoke suites. New model families
require an explicit scope change.

\section{Project artifacts}
The repository contains the detailed \href{project-report.md}{project report},
\href{article-podcast-plan.md}{article and podcast plan}, \href{bug-ledger.md}{bug ledger}, and
\href{coverage-matrix.md}{coverage matrix}. Component notes document the mapping from Parsl source
paths to TLA+ modules and runtime probes.

\end{document}
